\section{Introduction}
\label{sec:intro}

Can cheap data triage cartel investigations before costly
forensics? The empirical industrial-organization literature
reads collusion off the distribution of submitted bids
\citep{porter1993detection,bajari2003deciding,imhof2019detecting,wallimann2023machine,chassang2022robust}.
Yet most antitrust enforcers see only who won a public
contract and at what price---not what every bidder offered.
Electronic procurement systems archive the award envelope at
an analytical-warehouse interface that returns query results
in minutes; per-bidder bid amounts, where retained at all, sit
in operational logs reachable only after weeks of administrative
back-and-forth (\S\ref{sec:institutional}). Methods designed
for the bid layer simply do not run on the layer that
survives.

The enforcement-design literature
\citep{becker1968crime,stigler1970optimum,baker2002case,harrington2008detecting}
prescribes a response to this asymmetry: investigative
resources should be sequenced. A cheap screening stage triages
the costly forensic stage. The screening stage runs on the
data the agency already routinely observes; the forensic stage
inherits a smaller pool. We propose an enforcement architecture
along these lines, in which an award-layer screening stage
triages firms and environments before a costly bid-layer
forensic stage interrogates them. We instantiate the
architecture with a \emph{frequent-loser} flag built from
São Paulo's BEC contract-award records
($\valYearStart$--$\valYearEnd$, $\valSampleN$ tender-items).

A simple separating-equilibrium framework with cover bidders
and genuine entrants disciplines what the construct measures
(Online Appendix~A). Cover bidders in equilibrium satisfy
$\Pr(\text{win}\mid C)=0$: the $\text{wins}=0$ filter is a
separating choice, not a descriptive coincidence. Among
always-loser firms, $\log(1+\text{tenders\_count})$ ranks
firms by their posterior probability of being type $C$, and
endogenous sorting into cartel-affected items is part of what
the screening object measures rather than a confound to clean.
The binary $\text{median}+1.5\!\times\!\text{IQR}$ rule
coarsens this continuous primitive, and we call firms above
the threshold \emph{frequent losers}.

Three findings emerge from the empirical demonstration.
\textit{First}, the flag cuts the forensic stage's
bid-microdata pool by $\valArchFootprintReduction$
($\valArchSeqKone$ of $\valArchPoolSize$ firms) while still
flagging $\valArchTPThouSeqK$ of $\valArchNPos$ adjudicated
cartel cobidders. The architecture buys substantial
data-envelope economy without sacrificing recall against the
externally-validated cartel-adjacent population
(\S\ref{sec:forensic}). \textit{Second}, using only
contract-award data, the flag matches the seven-feature
Imhof--Wallimann bid-distribution pipeline that requires bid
microdata, and adds non-redundant signal when the two are
combined ($+\valImhofIncFL$ AUC, DeLong
$p=\valImhofIncFLp$). The screen and the bid-distribution
methods are informational complements, not substitutes
(\S\ref{sec:forensic}). \textit{Third}, built on
$2009$--$2016$ participation only, the screen prospectively
flags adjudicated cobidders in $2017$--$2019$ at firm-level
AUC $\valAUCFLfirmTemp$ under temporal holdout. The
discrimination is out-of-sample, not retrospective fitting
(\S\ref{sec:cade}).

The cobidder population behaves like the modeled type.
Cobidders deploy at higher intensity than FL non-cobidders
(Cohen's $d=\valBridgeUniqWinD$ on unique winners faced), bid
proximally to direct CADE defendants ($d=\valBridgeDirectD$),
and concentrate participation in narrower markets
($d=\valBridgeHHID$ on portfolio item-group HHI). At the bid
level they bid plausibly close to winners with elevated
within-firm cross-bid dispersion ($p<10^{-3}$ in a multivariate
logit holding participation intensity constant). The signature
matches credible cover bidding---bids low enough to be
plausible, dispersed enough to track rotation across cover
roles---rather than the textbook deliberately-uncompetitive
cover bid the early literature described
(\S\ref{sec:cade_bridge}). Frequent-loser presence is also
associated with a $\valHeadlineRange$ higher conditional log
negotiated unit price across four estimators on the broad
sample, with the positive sign concentrated in the
largest-tender-value stratum and a sign reversal under overlap
restriction; we report this descriptively
(\S\ref{sec:results}) and the paper does not rest on it.

The contribution sits at the intersection of four literatures.
Bid-coordination tests
\citep{bajari2003deciding,porter1993detection,conley2016detecting}
require a researcher-supplied partition between colluders and
a competitive fringe; the construct supplies a candidate
\emph{ex ante} firm-level flag computable from award records.
Bid-distribution screens
\citep{imhof2019detecting,wallimann2023machine,huber2019machine}
require bid microdata at the analytical-warehouse layer; the
construct operates on the operational layer and integrates
with bid-distribution methods at the forensic stage where
microdata are recoverable. Empirical cover-bidding work
identifies cover bidders \emph{ex post}, through leniency
testimony or judicial discovery
\citep{porter1999ohio,pesendorfer2000study,asker2010leniency,kawai2022detecting};
we are not aware of a published prospective identification of
cover-bidder candidates from public award-record data alone.
The enforcement-design literature
\citep{becker1968crime,stigler1970optimum,harrington2008detecting,baker2002case}
prescribes how detection technology and information
acquisition interact under costly observation; the
screening-then-forensic sequencing this paper proposes is a
particular structure of that information chain.

Four claims the paper does not make. The construct does not
adjudicate cartel membership: cobidders are the validation
object the data layer supports, and the AUC asymmetry against
direct CADE defendants is the design's empirical signature of
the loser-side scope, not a failure of the screening logic.
The pricing imprint is not a causal estimate of cover
bidding's effect on prices. The buyer-size gradient is scope
information about where the screening signal varies across
detection regimes, not identification of an institutional
channel. The preg\~ao--convite modal asymmetry is scope
information about where the screening object discriminates
better, not a positive test of the minimum-bidder-rule
mechanism.

\medskip\noindent
\textit{Roadmap.}
\S\ref{sec:literature} situates the contribution against the
bid-coordination, bid-distribution, cover-bidding, and
enforcement-design literatures.
\S\ref{sec:institutional} describes the BEC observability
layers and the CADE adjudication portfolio.
\S\S\ref{sec:data}--\ref{sec:empirical} state the
participation primitive, the binary rule, and the empirical
strategy.
\S\ref{sec:cade} carries the discrimination evidence and the
within-stratum theory-validation bridge.
\S\ref{sec:results} reports the price imprint as descriptive
corroboration and the segment-level decomposition of the sign
reversal.
\S\S\ref{sec:mechanisms}--\ref{sec:robustness} develop the
modal contrast, the detection-regime gradient, and the
threshold and identification audits.
\S\ref{sec:forensic} carries the architectural test against
bid-distribution detectors trained on bid microdata, including
the sequential-gatekeeper analysis and its temporal-holdout
audit. \S\S\ref{sec:limitations}--\ref{sec:conclusion} catalog
scope limitations and conclude.
